\subsection{Physical endpoint string order in the AKLT and cluster states}
\label{sec:physical_string_order_examples}

The following computations concern the stationary infinite-chain correlator
of Definition~\ref{def:physical_string_order_param}. They use physical
endpoint operators on the printed tensors, rather than arbitrary virtual
boundaries. The middle-string length is finite, but the boundary density
represents the thermodynamic state; these are not finite periodic-ring
expectations. They supply Examples~1--2 of
\cite[local TeX lines~392--411]{PerezGarcia2008StringOrder} without invoking
the stronger one-site injectivity hypotheses of the general virtual-boundary
theorems.

\begin{definition}[Physical data of the AKLT string correlator]
    \label{def:aklt_physical_string_data}
    \lean{MPSTensor.akltSpinRotationZ,
        MPSTensor.akltPGWSVC08Tensor,
        MPSTensor.akltPGWSVC08Tensor_apply,
        MPSTensor.akltPGWSVC08_rephasing_unitary,
        MPSTensor.akltPGWSVC08_sigmaZ_intertwine,
        MPSTensor.akltSpinRotationZ_eq_exp,
        MPSTensor.akltSpinRotationZ_unitary_ne_one}
    \leanok
    \uses{def:aklt_tensor, def:aklt_polynomial_hamiltonian}
    In the order $(m=0,+1,-1)$, the source tensor has letters
    \[
        B^0=\frac{\sigma_z}{\sqrt3},\qquad
        B^1=\sqrt{\frac23}\,\sigma_+,\qquad
        B^2=\sqrt{\frac23}\,\sigma_-.
    \]
    It is the diagonal physical rephasing $\operatorname{diag}(1,1,-1)$ of
    Definition~\ref{def:aklt_tensor}; this is a unitary basis change, not a
    virtual gauge. Take
    $x=y=S_z=\operatorname{diag}(0,1,-1)$ and
    $u=e^{i\pi S_z}=\operatorname{diag}(1,-1,-1)$.
    The rephasing leaves these physical operators unchanged, and the virtual
    gauge $V=\sigma_z$ implements $u$ on the printed letters.
\end{definition}

\begin{theorem}[The AKLT physical string-order value]
    \label{thm:aklt_physical_string_order_value}
    \lean{MPSTensor.twistedTransferMap_akltPGWSVC08_diagonal,
        MPSTensor.twistedTransferMap_aklt_diagonal,
        MPSTensor.akltPGWSVC08_canonical,
        MPSTensor.twistedTransferMap_aklt_spinZ_one,
        MPSTensor.twistedTransferMap_aklt_rotationZ_sigmaZ,
        MPSTensor.twistedTransferMap_aklt_spinZ_sigmaZ,
        MPSTensor.physicalStringOrderParam_aklt,
        MPSTensor.physicalStringOrderParam_akltPGWSVC08,
        MPSTensor.physicalStringOrderParam_akltPGWSVC08_identity_boundary,
        MPSTensor.akltPGWSVC08_hasPhysicalStringOrderWith,
        MPSTensor.akltPGWSVC08_hasPhysicalStringOrder}
    \leanok
    \uses{def:aklt_physical_string_data, def:physical_string_order_param,
        def:has_physical_string_order}
    The tensor $B$ is unital and has positive trace-one stationary density
    $\Lambda=\Id/2$. For every $N\geq0$,
    \begin{align}
        S_N(S_z,S_z,u) & =-\frac49.
        \label{eq:aklt_physical_string_order_value}
    \end{align}
    Thus the actual unblocked physical endpoints witness string order with
    limiting absolute value $4/9$ and a nonidentity physical unitary.
    The same value holds for the existing representative $A$ of
    Definition~\ref{def:aklt_tensor}.
\end{theorem}
\begin{proof}\leanok
    Diagonal insertions are invariant under the physical rephasing. Direct
    calculation gives
    \[
        \mathcal E_{S_z}(\Id)=\frac23\sigma_z,\qquad
        \mathcal E_u(\sigma_z)=\sigma_z,\qquad
        \mathcal E_{S_z}(\sigma_z)=-\frac23\Id.
    \]
    Iteration and the normalized trace give $-4/9$. The positive limit is
    constant, so no asymptotic approximation is needed.
\end{proof}

\begin{remark}[Normalization in the printed AKLT example]
    % **Local fix (normalized stationary density):** The printed identity
    % boundary conflicts with the source normalization. The formal value uses
    % I/2, with both readings proved and recorded in
    % docs/paper-gaps/pgwsvc08_string_order_virtual_boundary.tex.
    \textbf{Local fix:} The detailed correction is recorded in
    \path{docs/paper-gaps/pgwsvc08_string_order_virtual_boundary.tex}.
    The source's display \texttt{fixed} requires $\tr\Lambda=1$, but its
    AKLT example prints $\Lambda=\Id$. In dimension two the normalized
    density is $\Id/2$. Using the literal identity in the same correlator
    gives $-8/9$, as also proved above; the stated value $-4/9$ therefore
    requires this normalization correction.
\end{remark}

\begin{theorem}[The cluster state with one-site physical endpoints]
    \label{thm:cluster_one_site_physical_string_obstruction}
    \lean{MPSTensor.clusterTensorRMP_twistedTransferMap,
        MPSTensor.clusterTensorRMP_transferMap_one,
        MPSTensor.clusterTensorRMP_adjoint_fixes_maximallyMixed,
        MPSTensor.clusterTensorRMP_maximallyMixed_posDef_trace,
        MPSTensor.clusterTensorRMP_twistedTransferMap_pauliY,
        MPSTensor.clusterTensorRMP_pauliY_intertwine,
        MPSTensor.clusterTensorRMP_trace_pauliY_mul_letter_mul_conjTranspose,
        MPSTensor.clusterTensorRMP_twist_twice_endpoint,
        MPSTensor.clusterTensorRMP_physicalStringOrderParam_eq_zero,
        MPSTensor.clusterTensorRMP_not_hasPhysicalStringOrderWith}
    \leanok
    \uses{def:cluster_tensor_review, def:physical_string_order_param,
        def:has_physical_string_order}
    For the printed cluster tensor $A^0=\ket{0}\bra{+}$,
    $A^1=\ket{1}\bra{-}$, let $u=-\sigma_x$, $V=\sigma_y$ and
    $\Lambda=\Id/2$. The transfer map is unital, $\Lambda$ is a positive
    trace-one stationary density, and $V$ implements the local symmetry.
    Every pair of physical indices satisfies
    $\tr(V\Lambda A_nA_j^\dagger)=0$. For arbitrary one-site physical
    endpoints $x,y$,
    \[
        S_N(x,y,-\sigma_x)=0\qquad (N\geq2).
    \]
    Hence no such endpoints give a positive limiting absolute value for
    this fixed twist. This does not quantify over other possible twists.
\end{theorem}
\begin{proof}\leanok
    Every $\mathcal E_y(\Id)$ is diagonal. The twist sends $\Id$ to zero,
    $\sigma_z$ to $-\sigma_x$, and $\sigma_x$ to zero. Thus two iterations
    kill the right-endpoint image, and every later correlator vanishes.
\end{proof}

\begin{definition}[The cluster state's two-site physical endpoints]
    \label{def:cluster_two_site_string_endpoints}
    \lean{MPSTensor.clusterBlockedRMP_eq_mul,
        MPSTensor.clusterStringLeft,
        MPSTensor.clusterStringRight,
        MPSTensor.clusterStringTwist,
        MPSTensor.clusterStringLeft_apply_decode,
        MPSTensor.clusterStringRight_apply_decode,
        MPSTensor.clusterStringTwist_eq_blockKron}
    \leanok
    \uses{def:cluster_tensor_review}
    Take $x=\sigma_z\otimes\sigma_y$ and
    $y=\sigma_y\otimes\sigma_z$ on the two endpoint blocks. Blocked index
    $i$ decodes to the ordered word $(i\bmod2,\lfloor i/2\rfloor)$;
    the endpoint coefficient formulas use this order explicitly. On a
    two-site blocked chain the twist is $u\otimes u$.
\end{definition}

\begin{theorem}[The cluster physical string-order value]
    \label{thm:cluster_two_site_physical_string_order_value}
    \lean{MPSTensor.clusterBlockedRMP_twistedTransferMap_right,
        MPSTensor.clusterBlockedRMP_twistedTransferMap_left,
        MPSTensor.clusterTensorRMP_twoSite_physicalStringOrderParam,
        MPSTensor.clusterBlockedRMP_twistedTransferMap_pauliY,
        MPSTensor.clusterBlockedRMP_physicalStringOrderParam,
        MPSTensor.clusterBlockedRMP_hasPhysicalStringOrderWith,
        MPSTensor.clusterBlockedRMP_hasPhysicalStringOrder}
    \leanok
    \uses{def:cluster_two_site_string_endpoints,
        thm:cluster_one_site_physical_string_obstruction,
        def:has_physical_string_order}
    With $\mathcal E_x^{(2)},\mathcal E_y^{(2)}$ the insertions on the two-site
    endpoint blocks, every single-site middle length $N\geq0$ satisfies
    \begin{align}
        \tr\!\left[\frac{\Id}{2}\,
            \mathcal E_x^{(2)}\mathcal E_{-\sigma_x}^{N}
        \mathcal E_y^{(2)}(\Id)\right] & =1.
        \label{eq:cluster_two_site_physical_string_order_value}
    \end{align}
    Odd middle lengths are included. The two-site blocked tensor also has
    physical string-order parameter one and satisfies the positive-limit
    predicate with the nonidentity unitary $u\otimes u$.
\end{theorem}
\begin{proof}\leanok
    The explicit endpoint maps satisfy
    $\mathcal E_y^{(2)}(\Id)=-\sigma_y$ and
    $\mathcal E_x^{(2)}(\sigma_y)=-\Id$, while
    $\mathcal E_{-\sigma_x}(\sigma_y)=\sigma_y$. Both minus signs cancel,
    and the normalized trace is one. The same fixed-point calculation applies
    to the blocked twist.
\end{proof}
